\documentclass[main.tex]{subfiles}


\begin{document}

%from pagenumber to up
% \phantomsection 
% \hypertarget{toc}{}

 

 

 
   

\section{Потенциал}


В общем случае электрическое поле\footnote{Электрическое поле далее предполагается \emph{постоянным}~--- то есть \emph{электростатическим}.} способно разгонять заряд, помещенный в него,~--- говорят, что заряд в таком поле имеет потенциальную энергию.

\sbox{0}{%
    \begin{tikzpicture}[font=\small,            line cap=round,                             >={Stealth[inset=0pt,round,length=3mm, angle'=20]},vec/.style={>={Stealth[inset=0pt,round,length=3.5mm, angle'=20]}}]


%line
\foreach \i in{0,...,1}
\draw[Green,decoration={markings,mark=at position {1/2+1.5/20} with {\arrow{>}}},postaction={decorate}] ({\i*2},2) node [above,black] {$E$} --++ ({-90}:2);

%plane
\draw[dashed] (-.5,0) --++ (3,0) node[right,inner ysep=0] {П};


%size
\draw[dashed] (1.15,1) --++ (.4,0);
\draw[<->] (1.35,1) --++ (0,-1) node[right,midway] {$h$};


%charge
\shade[ball color=red] (1,1) circle (.15);
\node[above] at (1,1.15) {$q$};


    \end{tikzpicture}%
}

{%
\setlength\intextsep{0pt}
\begin{wrapfigure}[]{r}{\wd0}
    \centering
    \usebox{0}%
    \caption{Заряд в однородном поле}
    \label{fig:p101}
\end{wrapfigure}


Пусть заряд находится в однородном поле (рис.\,\ref{fig:p101}).

%     \begin{figure}[H]
%     \centering

%     \begin{tikzpicture}[font=\small,            line cap=round,                             >={Stealth[inset=0pt,round,length=3mm, angle'=20]},vec/.style={>={Stealth[inset=0pt,round,length=3.5mm, angle'=20]}}]


% %line
% \foreach \i in{0,...,1}
% \draw[Green,decoration={markings,mark=at position {1/2+1.5/20} with {\arrow{>}}},postaction={decorate}] ({\i*2},2) node [above,black] {$E$} --++ ({-90}:2);

% %plane
% \draw[dashed] (-.25,0) --++ (2.5,0) node[right] {П};


% %size
% \draw[dashed] (1.15,1) --++ (.4,0);
% \draw[<->] (1.35,1) --++ (0,-1) node[right,midway] {$h$};


% %charge
% \shade[ball color=red] (1,1) circle (.15);
% \node[above] at (1,1.15) {$q$};


%     \end{tikzpicture}
% \caption{Заряд в однородном поле}
% \label{fig:p101}
% \end{figure}



Положительный заряд вначале покоится и взаимодействует только с электрическим полем. Ясно, что заряд будет <<падать>> под действием этого поля вдоль его линий на воображаемую поверхность~П (поле может быть направлено в любую сторону, поэтому <<падать>> заряд может также в любую сторону). Тогда говорят, что заряд~$q$, находящийся в однородном электрическом поле~$E$ и способный <<упасть>> с высоты~$h$ под действием этого поля, обладает \textbf{потенциальной энергией} (обусловленной взаимодействием заряда с полем):
\begin{equation}
    \label{8poten_e1}
    W=qEh.
\end{equation}

}

Случай заряда в однородном электрическом поле напоминает ситуацию с телом, поднятым над поверхностью планеты. Как известно, потенциальная энергия тела, обусловленная притяжением тела и планеты, равна~$E_\text{п}=mgh$. Таким образом, в формуле~\eqref{8poten_e1} заряд~$q$ является аналогом массы~$m$, а напряженность~$E$ является аналогом ускорения свободного падения~$g$. 

\textbf{Потенциал} ($\varphi$ [В])~--- это энергетическая характеристика электрического поля в точке пространства:
\begin{equation}
    \varphi=\frac{W}{q_\text{внесен}},
\end{equation}
где $W$~--- потенциальная энергия заряда~$q_\text{внесен}$ в электрическом поле.

% Потенциал численно равен энергии, которую имел бы заряд 1~Кл в данном электрическом поле в рассматриваемой точке.

\textbf{Потенциал поля точечного заряда} в вакууме
% \footnote{В диэлектрике формула~\eqref{e31} приобретает вид: $\varphi=k\dfrac{q_\text{источника}}{\varepsilon r}$.}
находят по формуле: 
\begin{equation}
    \label{8poten_e2}
    \varphi=k\frac{q_\text{источника}}{r},
\end{equation}
где $r$~--- расстояние от заряда~$q_\text{источника}$ до рассматриваемой точки.

В диэлектрике формула~\eqref{8poten_e2} приобретает вид: $\varphi=k\dfrac{q_\text{источника}\strut}{\varepsilon r\strut}$.

Если заряд перемещается в поле из точки~1 в точку~2, то полезно ввести следующую величину. \textbf{Напряжение} ($U$ [В])~--- это \emph{разность потенциалов}:
\begin{equation}
    U=\varphi_1-\varphi_2=\frac{A}{q_\text{внесен}},
\end{equation}
где $\varphi_1$ и $\varphi_2$~--- потенциалы поля в начальной и конечной точке траектории, $A$~--- работа силы поля, действующей на перемещающийся заряд~$q_\text{внесен}$ ($A=-\Delta W$).

Важно подчеркнуть, что \emph{разность потенциалов есть потенциал начальной точки минус потенциал конечной точки.}

Следует отметить, что \emph{вектор напряженности поля указывает направление убывания потенциала.} Так, на рис.\,\ref{fig:p101} в точке нахождения заряда~$q$ потенциал поля больше, чем потенциал точки поверхности~П под зарядом.

\textbf{Напряжение в однородном поле} равно:
\begin{equation}
    U=Ed,
\end{equation}
где $d$~--- расстояние между двумя точками \emph{вдоль линии поля}.


 
\end{document}